\documentclass[a4paper,11pt]{article}

\usepackage{revy}
\usepackage[utf8]{inputenc}
\usepackage[T1]{fontenc}
\usepackage[danish]{babel}


\revyname{MatematikRevyen}
\revyyear{2026}
\version{0.1}
\eta{$1.5$ minutter}
\status{Færdig}

\title{Vektorrum}
\author{Louie '24 og Christian '24}
\melody{Skrillex: ``Bangarang''}

\begin{document}
\maketitle
\enlargethispage{3\baselineskip}

\begin{roles}
\role{S}[Skuespiller] Sanger
\role{Dansere}[Dansere] Vektorer, kommer ind i omkvædet
\end{roles}

\begin{props}
\prop{Pile}[Person, der skaffer] Tre pile til hver danser: én langs hver arm og én op ad kroppen
\end{props}

\scene{S står alene på scenen. I omkvædet kommer danserne ind med pile langs kroppen.}

\begin{song}
  \sings{S}[vers] Additiv med invers
  Addi-addi-addi-addi-addi-addiv med invers
  og lukket
  distributiv med enhed
  d-d-d-d-d-d distrubutiv med enhed
  og abelsk
  kommutativ ved addition
  k-k-k-k-k-kommutativ ved addition
  og skalar
  associativ med enhed
  ass-ass-ass-ass-associativ med enhed

  {\it (Danserne kommer ind med pile langs kroppen)}

  \sings{S}[omkv] Vektorrum
  Vektorrum
  Vektorrum
  Vektorrum
  Vektorrum
  Vektorrum
  Vektorrum
  Vektorrum
\end{song}

\end{document}
